\documentclass[11pt,oneside,openany]{book}

\usepackage[T1]{fontenc}
\usepackage{lmodern}
\usepackage[letterpaper,margin=1.10in,includeheadfoot]{geometry}
\usepackage{amsmath,amssymb,amsthm,mathtools}
\usepackage{microtype}
\usepackage{xcolor}
\usepackage{titlesec}
\usepackage{fancyhdr}
\usepackage{needspace}
\usepackage{hyperref}

\ifdefined\pdftrailerid
  \pdftrailerid{MathGenome ALG-02 core supplement}
\fi

\definecolor{ink}{HTML}{172923}
\definecolor{deepgreen}{HTML}{2B5242}
\definecolor{rulegreen}{HTML}{8A9F92}
\definecolor{palegreen}{HTML}{F0F4EF}
\hypersetup{
  colorlinks=true,linkcolor=deepgreen,urlcolor=deepgreen,citecolor=deepgreen,
  pdftitle={Group Theory},
  pdfsubject={Chief factors, extensions and group cohomology, and profinite completion},
  pdfauthor={MathGenome},bookmarksnumbered=true
}
\setcounter{tocdepth}{1}
\setcounter{secnumdepth}{2}
\setlength{\headheight}{15pt}
\setlength{\emergencystretch}{2em}
\titleformat{\chapter}[display]
  {\normalfont\color{ink}}{\color{deepgreen}\large\scshape Chapter \thechapter}{9pt}
  {\huge\bfseries}
\titlespacing*{\chapter}{0pt}{4pt}{24pt}
\titleformat{\section}{\normalfont\Large\bfseries\color{ink}}{\thesection}{0.7em}{}
\titleformat{\subsection}{\normalfont\large\bfseries\color{ink}}{\thesubsection}{0.7em}{}

\pagestyle{fancy}
\fancyhf{}
\fancyhead[L]{\small\sffamily\color{deepgreen}MathGenome}
\fancyhead[R]{\small\nouppercase{\leftmark}}
\fancyfoot[C]{\small\thepage}
\renewcommand{\headrulewidth}{0.35pt}
\renewcommand{\chaptermark}[1]{\markboth{\thechapter.\ #1}{}}
\fancypagestyle{plain}{%
  \fancyhf{}\fancyfoot[C]{\small\thepage}\renewcommand{\headrulewidth}{0pt}}

\newtheorem{theorem}{Theorem}[chapter]
\newtheorem{lemma}[theorem]{Lemma}
\newtheorem{proposition}[theorem]{Proposition}
\newtheorem{corollary}[theorem]{Corollary}
\theoremstyle{definition}
\newtheorem{definition}[theorem]{Definition}
\newtheorem{example}[theorem]{Example}
\theoremstyle{remark}
\newtheorem{remark}[theorem]{Remark}
\allowdisplaybreaks[1]

\begin{document}
\frontmatter
\hypersetup{pageanchor=false}
\begin{titlepage}
  \thispagestyle{empty}
  \thispdfpagelabel{Cover}
  \begin{flushleft}
    \vspace*{0.65in}
    {\large\scshape\color{deepgreen}MathGenome\par}
    \vspace{0.32in}
    {\color{rulegreen}\rule{0.94\textwidth}{0.65pt}\par}
    \vspace{0.45in}
    {\color{ink}\fontsize{34}{39}\selectfont\bfseries
       Group Theory\par}
    \vspace{0.30in}
    {\color{deepgreen}\Large
       Chief Series, Extensions,\\[3pt] and Profinite Completion\par}
    \vfill
    {\large Supplementary lecture notes\par}
    \vspace{0.12in}
    {\normalsize ALG-02 \enspace\textperiodcentered\enspace MathGenome v0.22.0 reference baseline\par}
    \vspace*{0.48in}
  \end{flushleft}
\end{titlepage}
\hypersetup{pageanchor=true}
{\small\tableofcontents}
\chapter*{Scope and reading guide}
\addcontentsline{toc}{chapter}{Scope and reading guide}
\markboth{Scope and reading guide}{}

These notes supplement the group theory in \emph{Abstract Algebra}
with finite chief series, extension classes, and profinite completion.
The guide traces the six core modules of ALG-02 through the existing
proofs and the three chapters of this volume.

{\setlength{\fboxsep}{0.8em}\setlength{\fboxrule}{0.35pt}%
\noindent\fcolorbox{rulegreen}{palegreen}{%
\begin{minipage}{\dimexpr\linewidth-2\fboxsep-2\fboxrule\relax}
\textbf{Result keys.} \texttt{GR02-C02.01} belongs to module
\texttt{ALG-02-C02}. The source map records the PDF and \TeX\
locations and proof dependencies. Keys beginning with \texttt{K07}
refer to the released D07 notes.
\end{minipage}}\par}

\section*{The six core modules}

\paragraph{C01 \enspace Groups, quotients, and actions.}
D07, groups K07.01--02 and K07.04--05, develops generators, cyclic
and permutation groups, quotient groups and the isomorphism theorems,
and orbit--stabilizer and counting arguments \cite{d07}.

\paragraph{C02 \enspace Finite groups and series.}
Cauchy, Sylow, and composition series are in K07.06, solvability
in K07.15, and nilpotence in K07.18. Chapter~\ref{ch:gr02-chief} adds
chief factors and their comparison as groups with an ambient action.

\paragraph{C03 \enspace Abelian and free groups.}
D07, K07.07, proves the finitely generated abelian structure theorem.
K07.01 gives free groups and presentations. For Nielsen--Schreier
and the finite-index rank formula, see \cite[Results B04--B05]{geomgroup}.

\paragraph{C04 \enspace Extensions and group cohomology.}
D07, K07.04, treats the semidirect product and the splitting test.
See \cite[Results A02--A03, A05, and A10]{homalg} for the bar
constructions and the cohomological five-term sequence.
Chapter~\ref{ch:gr02-extensions} develops classification,
Schur multipliers, and universal central extensions.

\paragraph{C05 \enspace Finite quotients and inverse limits.}
For the Galois constructions and the finite-field case
\(\operatorname{Gal}(\overline{\mathbb F}_q/\mathbb F_q)
\cong\widehat{\mathbb Z}\), see \cite[Results T04--T06 and E06]{field}.
Chapter~\ref{ch:gr02-profinite} constructs the completion of an
arbitrary abstract group.

\paragraph{C06 \enspace Combinatorial and geometric methods.}
\emph{Geometric Group Theory} covers Cayley and Schreier graphs, word
problems, tree actions, and growth, with the hypotheses of each
algorithm and comparison spelled out \cite[Chs.~1--3 and~6]{geomgroup}.

\mainmatter
\chapter{Chief series of finite groups}
\label{ch:gr02-chief}

The composition series in D07 have simple factors and require
normality of each term in its successor
\cite[entries K07.06.composition\_series\_definition and
K07.06.jordan\_holder\_extension]{d07}. Here we refine the normal
subgroups of a fixed ambient group \(G\). Their successive quotients
carry the conjugation action of \(G\), which the factor comparison
will preserve.

\section{Chief factors and conjugation}

\begin{definition}[GR02-C02.01. Chief series and chief factors]
\label{def:gr02-chief-series}
Let \(G\) be a finite group. A chief series is a strictly increasing
chain
\[
  1=N_0<N_1<\cdots<N_r=G
\]
of normal subgroups of \(G\) such that no normal subgroup of \(G\)
lies strictly between \(N_{i-1}\) and \(N_i\), for \(1\le i\le r\).
The quotients \(N_i/N_{i-1}\) are the chief factors, and the integer
\(r\) is the length. The chain of length zero is allowed for \(G=1\).

A \(G\)-group is a group equipped with an action of \(G\) by
automorphisms. When \(K\le L\) are both normal in \(G\), the
natural \(G\)-group structure on \(L/K\) is
\[
  g\cdot(xK)=(gxg^{-1})K \qquad (g\in G,\ x\in L).
\]
A \(G\)-isomorphism of two such factors is a group isomorphism
which intertwines the two actions of the same ambient \(G\).
\end{definition}

The formula is well defined since conjugation by \(g\) preserves both
\(K\) and \(L\). The action restricted to \(L\) induces every inner
automorphism of \(L/K\), so a \(G\)-invariant subgroup is normal in
\(L/K\). For the trivial group, the list of factors is empty.

\begin{proposition}[GR02-C02.02. Existence and the minimal normal criterion]
\label{prop:gr02-chief-existence}
Every finite group has a chief series. A strictly increasing chain as
in Definition~\ref{def:gr02-chief-series}, with all its terms normal in
\(G\), is a chief series precisely when \(N_i/N_{i-1}\) is a minimal
nontrivial normal subgroup of \(G/N_{i-1}\) for each \(i\).
\end{proposition}
\begin{proof}
The finite poset of normal subgroups of \(G\) has a maximal chain
from \(1\) to \(G\). Such a chain admits no refinement by a normal
subgroup of \(G\). Normal subgroups of \(G/N_{i-1}\) contained in
\(N_i/N_{i-1}\) correspond under the quotient map to normal
subgroups of \(G\) between \(N_{i-1}\) and \(N_i\)
\cite[entry K07.04.correspondence\_theorem]{d07}. The endpoints of
the interval correspond to the trivial subgroup and the whole
factor. This proves the criterion in both directions.
\end{proof}

\section{Uniqueness and the solvable case}

\begin{theorem}[GR02-C02.03. Equivariant Jordan--H\"older uniqueness for chief factors]
\label{thm:gr02-chief-uniqueness}
Any two chief series of a finite group \(G\) have the same length
and the same multiset of \(G\)-isomorphism classes of factors.
\end{theorem}
\begin{proof}
Keep \(G\) fixed throughout. For every normal subgroup \(N\) of
\(G\), consider saturated chains from \(1\) to \(N\) in the poset
of normal subgroups of \(G\). Here saturated means that consecutive
terms admit no intermediate term in this poset. Such chains exist
by finiteness, and their factors all carry the \(G\)-action of
Definition~\ref{def:gr02-chief-series}.

We compare the multisets of factors of these chains by induction on
\(|N|\). The case \(N=1\) has an empty multiset. Let \(A\) and \(B\)
be the penultimate terms of two chains ending at a nontrivial
\(N\). If \(A=B\), induction compares the two chains up to \(A\)
and the top factor \(N/A\) is common to both. If \(A\ne B\), neither
contains the other by maximality. The product \(AB\) is normal in
\(G\) and strictly contains \(A\), so \(AB=N\). Put \(L=A\cap B\).
The quotient maps give isomorphisms
\[
  A/L \longrightarrow N/B,\quad aL\longmapsto aB,
  \qquad
  B/L \longrightarrow N/A,\quad bL\longmapsto bA.
\]
In each case the indicated intersection is the kernel before
passing to the quotient, and \(AB=N\) gives surjectivity.
Conjugating representatives by any \(g\in G\) commutes with both
maps, so these are \(G\)-isomorphisms.

There is no proper nontrivial \(G\)-invariant subgroup of \(N/B\),
since its inverse image in \(N\) would be normal in \(G\) and
strictly between \(B\) and \(N\). The same holds for \(N/A\).
By the two isomorphisms, \(L<A\) and \(L<B\) are therefore saturated
steps in the poset. Extend one saturated chain from \(1\) to \(L\)
first by \(A\) and then by \(N\), or first by \(B\) and then by \(N\).
Induction within \(A\) and \(B\) compares the original chains with
these two chains, including the common top step on each side.
The two constructed chains share all factors through \(L\), and
their last two factors match in reverse order by the displayed
\(G\)-isomorphisms. This finishes the induction. Taking \(N=G\)
proves the theorem.
\end{proof}

This is the descending comparison used for ordinary composition
series in D07 with the ambient conjugation retained at every step
\cite[entry K07.06.jordan\_holder\_extension]{d07}. The ordinary
group-isomorphism comparison for chief factors is also stated in
\cite[Theorem~5.9]{chief-source}.

\begin{corollary}[GR02-C02.04. Chief factors of finite solvable groups]
\label{cor:gr02-chief-solvable}
A finite group \(G\) is solvable if and only if it has a chief
series with elementary abelian factors. If \(G\) is solvable, every
chief factor is the additive group of a nonzero finite-dimensional
\(\mathbb F_p\)-vector space for some prime \(p\) depending on the
factor, and its \(G\)-action is linear over that field.
\end{corollary}
\begin{proof}
Assume \(G\) solvable. Passage to subgroups and quotients preserves
solvability by \cite[entry K07.15.solvable\_closure]{d07}. Hence the
chief factor \(D=N_i/N_{i-1}\) is solvable. Its derived subgroup
\([D,D]\) is characteristic in \(D\), hence is preserved by
conjugation by \(G\). Since \(D\) is nontrivial and solvable,
\([D,D]\) is proper. Minimality forces \([D,D]=1\), so \(D\) is
abelian. Choose a prime \(p\) dividing the order of some nonidentity
element of \(D\). A power of that element has order \(p\), and
\[
  D[p]=\{x\in D : x^p=1\}
\]
is a nontrivial characteristic subgroup because \(D\) is abelian.
Minimality now gives \(D[p]=D\). An abelian group of exponent \(p\)
has its usual \(\mathbb F_p\)-vector space structure by repeated
addition, and each group automorphism respects this scalar action.

Conversely, if the factors of a chief series are elementary
abelian then they are solvable. Closure under extensions makes
each term in the series solvable, starting at \(1\). This includes
the chain of length zero for \(G=1\).
\end{proof}

\begin{example}[GR02-C02.05. The two kinds of series in the alternating group]
\label{ex:gr02-chief-a4}
In \(A_4\), let
\[
  V=\{1,(12)(34),(13)(24),(14)(23)\},\qquad
  U=\langle(12)(34)\rangle .
\]
The subgroup \(V\) is elementary abelian of order \(4\), since the
product of two distinct double transpositions in \(V\) is the
third. It is normal in \(A_4\) since conjugation preserves cycle
type. Conjugation by \((123)\) permutes the three nonidentity
elements of \(V\) cyclically. Consequently every nontrivial
normal subgroup of \(A_4\) contained in \(V\) equals \(V\).
As \(A_4/V\) has order \(3\), \(1<V<A_4\) is a chief series.
Its factors are \(V\) and \(C_3\). The action on \(V\) factors
through \(A_4/V\) and permutes its three nonzero vectors, while
the conjugation action on the factor \(A_4/V\) is trivial.

The longer chain \(1<U<V<A_4\) is a composition series since
adjacent terms are normal and the factors have prime orders
\(2,2,3\). The subgroup \(U\) is moved to another subgroup by
\((123)\), so this chain fails the ambient normality condition
in Definition~\ref{def:gr02-chief-series}. See also
\cite[Examples~5.2 and 5.7]{chief-source}.
\end{example}
\chapter{Group extensions and the Schur multiplier}
\label{ch:gr02-extensions}

The semidirect product and its internal criterion are established
in \cite[K07.04.semidirect\_product and K07.04.semidirect\_recognition]{d07}.
Here an extension of \(Q\) by \(N\) comes with a specified
inclusion \(N\hookrightarrow E\) and quotient map \(p:E\to Q\).
We identify \(N\) with its image. A splitting is a homomorphism
\(s:Q\to E\) with \(ps=\operatorname{id}_Q\). Its image \(H\) is a
complement, meaning \(NH=E\) and \(N\cap H=\{1\}\). Conversely the
inverse of \(p|_H\) is a splitting for every complement. All groups
in this chapter are discrete, and the abelian kernel \(A\) in a
cohomology group is written additively. Integral group homology
uses the trivial coefficient action unless otherwise specified.

\section{Factor systems and outer actions}

\begin{proposition}[GR02-C04.01. Factor systems for a general kernel]
\label{prop:gr02-factor-system}
Choose a set-theoretic section \(s:Q\to E\) with \(s(1)=1\).
Define \(\alpha_q\in\operatorname{Aut}(N)\) by conjugation by \(s(q)\),
and put \(f(q,r)=s(q)s(r)s(qr)^{-1}\in N\). Then
\begin{align*}
 \alpha_q\alpha_r &=\operatorname{Inn}(f(q,r))\,\alpha_{qr},\\
 f(q,r)f(qr,t)&=\alpha_q(f(r,t))f(q,rt).
\end{align*}
Here \(\operatorname{Inn}(n)(x)=nxn^{-1}\). If \(s'(q)=b(q)s(q)\)
with \(b(1)=1\), the new data are
\[
 \alpha'_q=\operatorname{Inn}(b(q))\alpha_q,\qquad
 f'(q,r)=b(q)\alpha_q(b(r))f(q,r)b(qr)^{-1}.
\]
In particular \(q\mapsto[\alpha_q]\) is a homomorphism
to \(\operatorname{Out}(N)=\operatorname{Aut}(N)/\operatorname{Inn}(N)\)
independent of the section. Conversely, data satisfying the two
identities and the normalizations
\(\alpha_1=\operatorname{id}_N\), \(f(q,1)=f(1,q)=1\)
give an extension on \(N\times Q\) with product
\[
 (n,q)(m,r)=\bigl(n\,\alpha_q(m)f(q,r),qr\bigr).
\]
\end{proposition}

\begin{proof}
Conjugation by \(s(q)s(r)=f(q,r)s(qr)\) gives the first identity.
Expanding \((s(q)s(r))s(t)=s(q)(s(r)s(t))\) and cancelling \(s(qrt)\)
on the right gives the second. Substitution of \(b(q)s(q)\) gives
the displayed change of section. The inner automorphisms form a
normal subgroup of \(\operatorname{Aut}(N)\), since conjugating
\(\operatorname{Inn}(n)\) by any automorphism \(\beta\) gives
\(\operatorname{Inn}(\beta(n))\). Thus the first identity and the
change of section give the outer homomorphism.

For the converse, if \(u\in N\), the two associations of
\((n,q)(m,r)(u,t)\) have first coordinates
\[
 n\alpha_q(m) f(q,r)\alpha_{qr}(u) f(qr,t)
 \quad\text{and}\quad
 n\alpha_q(m)\alpha_q\alpha_r(u)\alpha_q(f(r,t))f(q,rt).
\]
They agree by the two identities, and the second coordinates both
equal \(qrt\). The normalizations give the identity \((1,1)\).
A right inverse to \((n,q)\) can be found with second coordinate
\(q^{-1}\) by solving \(\alpha_q(m)=n^{-1}f(q,q^{-1})^{-1}\).
A left inverse with that second coordinate is obtained by solving
\(m\alpha_{q^{-1}}(n)f(q^{-1},q)=1\).
Associativity makes the two inverses equal. Projection is a
surjective homomorphism with kernel \(N\times\{1\}\).
For the original extension the coordinate map
\((n,q)\mapsto ns(q)\) is bijective and respects the product.
\end{proof}

When \(N=A\) is abelian, \(\alpha_q\) is independent of the lift of
\(q\) and is an ordinary action of \(Q\) on \(A\). An equivalence
between extensions with fixed kernel and quotient is a homomorphism
of middle groups inducing the identity on \(A\) and on \(Q\).
Such a map is injective because its kernel lies in \(A\). For
surjectivity, first lift the quotient coordinate of a target
element and then adjust the lift by an element of \(A\).

\section{Abelian kernels and low-degree cohomology}

We use the inhomogeneous bar differential from
\cite[A02 and A03]{homalg}, with the left \(Q\)-module structure on
\(A\). For \(a:Q\to A\), the two-coboundary is
\[
 (\delta a)(q,r)=a(q)+q\cdot a(r)-a(qr).
\]
A two-cocycle \(c\) satisfies
\[
 c(q,r)+c(qr,t)=q\cdot c(r,t)+c(q,rt).
\]
It is normalized if \(c(1,q)=c(q,1)=0\). The normalized
cocycles give the same \(H^2\) classes as all inhomogeneous
cocycles. Indeed the identity with \(q=1\) shows
\(c(1,r)=c(1,1)\), and with \(t=1\) shows \(c(q,1)=q\cdot c(1,1)\).
Subtract the coboundary of the constant cochain with value
\(c(1,1)\) to normalize. If two normalized cocycles differ by
\(\delta a\), evaluating at \((1,r)\) gives \(a(1)=0\).

\begin{theorem}[GR02-C04.02. Abelian-kernel extensions and \(H^2\)]
\label{thm:gr02-h2-extensions}
Fix a left \(Q\)-action on \(A\). Equivalence classes of extensions
of \(Q\) by \(A\) inducing that action correspond bijectively to
\(H^2(Q,A)\). The zero class corresponds to the split extensions.
\end{theorem}

\begin{proof}
By Proposition~\ref{prop:gr02-factor-system}, a normalized
section gives a normalized cocycle \(c\), and changing the section
adds a coboundary. An extension equivalence preserves the cocycle
of the transported section. Conversely, a normalized cocycle
defines an extension by the product
\[
 (x,q)(y,r)=(x+q\cdot y+c(q,r),qr)
\]
on \(A\times Q\), by the converse in the same proposition.
Its canonical section \(q\mapsto(0,q)\) recovers \(c\), and the
conjugation action on \(A\) is the given action. If
\(c'=c+\delta a\), then \((x,q)\mapsto(x+a(q),q)\) carries the
product for \(c'\) to the product for \(c\) and fixes the two
end groups. Hence the constructions are inverse on classes.

A homomorphic section has cocycle zero. If the class is zero,
choose \(a\) with \(c+\delta a=0\); after the corresponding change
of section its cocycle is zero, so the section is homomorphic.
By the internal criterion this is precisely the semidirect product
for the fixed action.
\end{proof}

\begin{theorem}[GR02-C04.03. Complements and \(H^1\)]
\label{thm:gr02-h1-complements}
For a split extension with abelian kernel \(A\), choose a reference
splitting \(s_0\). The \(A\)-conjugacy classes of complements to
\(A\) in \(E\) are in bijection with \(H^1(Q,A)\), with the class
of \(s_0(Q)\) corresponding to zero.
\end{theorem}

\begin{proof}
Identify \(E=A\rtimes Q\) by \(s_0\). A section is the graph
of a function \(z:Q\to A\), and the semidirect product rule
makes it a homomorphism precisely when
\(z(qr)=z(q)+q\cdot z(r)\); this also implies \(z(1)=0\).
Every complement determines its unique splitting by projection.
Conjugating the graph by \((a,1)\) changes \(z(q)\) to
\(z(q)+a-q\cdot a\). As \(a\) varies these changes are exactly
the one-coboundaries, since they are the negatives of
\(q\mapsto q\cdot a-a\). The description of \(H^1\) in
\cite[A03]{homalg} now gives the bijection.
\end{proof}

The extension correspondence and the choice of reference
splitting are also treated in \cite{conrad-extensions} and
\cite[Theorems 23.3 and 24.3]{lassueur}. The cohomological
Lyndon--Hochschild--Serre sequence is proved in
\cite[A05]{homalg}. The next section constructs the low-degree
homological sequence directly on bar chains.

\begin{example}[GR02-C04.04. Two extensions of \(C_2\) by \(C_2\)]
\label{ex:gr02-c2-extensions}
Give \(A=C_2\) the trivial action of \(Q=C_2=\{1,t\}\).
A normalized cocycle is determined by \(c(t,t)\), which can
take either value in \(C_2\). Indeed, the identity holds whenever
an argument is \(1\), and the triple \((t,t,t)\) gives just
\(c(t,t)=c(t,t)\). Normalized coboundaries have
\((\delta a)(t,t)=2a(t)=0\). The class \(c(t,t)=0\) gives
\(C_2\times C_2\). In the other class the lift \((0,t)\) has
square \((1,1_Q)\), where \(1\) in the first coordinate denotes
the nonzero element of \(A=\mathbb Z/2\mathbb Z\).
It has order four and generates the middle group \(C_4\).
Thus even extensions with trivial action on the kernel can
represent different cohomology classes.
\end{example}

\section{The Hopf formula}

For groups \(U,V\) contained in a group \(E\), the notation
\([U,V]\) denotes the subgroup generated by
\([u,v]=uvu^{-1}v^{-1}\) for \(u\in U\) and \(v\in V\).
In particular \([E,E]\) is the commutator subgroup. The integral
bar complex from \cite[A10]{homalg}, on trivial coefficients,
identifies \(H_1(K;\mathbb Z)\) with the abelianization of \(K\):
in degree one the generators \([k]\) are cycles and the
boundaries are \([l]-[kl]+[k]\). They impose exactly the relation
that \(k\mapsto[k]\) is a homomorphism to an abelian group.

\begin{proposition}[GR02-C04.05. The homological five-term sequence]
\label{prop:gr02-homological-five-term}
For an extension \(1\to N\to E\xrightarrow{p}Q\to 1\), there is
a natural exact sequence
\[
 H_2(E;\mathbb Z)\longrightarrow H_2(Q;\mathbb Z)
 \xrightarrow{d^2_{2,0}}N/[E,N]\longrightarrow
 H_1(E;\mathbb Z)\longrightarrow H_1(Q;\mathbb Z)\longrightarrow 0.
\]
The maps between the homologies of \(E\) and \(Q\) are induced by
\(p\). The map from \(N/[E,N]\) is induced by inclusion.
\end{proposition}

\begin{proof}
Let \(C_\bullet(G)\) denote the integral normalized bar chains
with trivial coefficients, with any tuple containing \(1\) taken
as zero. The unnormalized integral bar chains compute group homology
by \cite[A10]{homalg}. To check normalization in the degrees at issue,
let \(D_\bullet\) be the subcomplex spanned by tuples containing \(1\).
It is a subcomplex because the two faces removing a chosen identity
cancel, and the other faces still contain it. Here \(D_0=0\) and
\(d[1\mid1]=[1]\) kills \(H_1(D)\). A generator of \(D_2\) has
boundary \([1]\). Modulo boundaries from \(D_3\) all generators of
\(D_2\) equal \([1\mid1]\), since
\[
 d[1\mid1\mid g]=[1\mid g]-[1\mid1],
 \qquad
 d[g\mid1\mid1]=[1\mid1]-[g\mid1].
\]
A cycle in \(D_2\) has coefficient sum zero, so \(H_2(D)=0\).
The quotient by \(D_\bullet\) thus has the same homology through
degree two. Projection is surjective on normalized chains in every degree. Put
\(J_\bullet=\ker(C_\bullet(E)\to C_\bullet(Q))\). Since \(J_0=0\),
its first homology is \(J_1/dJ_2\).

The group \(J_1\) is generated by differences \([nx]-[x]\), where
\(n\in N\) and \(x\in E\). Indeed, choose a representative \(s(q)\)
for each fiber, with \(s(1)=1\). Then \([e]-[s(p(e))]\) is a basis
of \(J_1\) as \(e\) ranges over the elements other than the chosen
representatives. Sending this basis element to the class of
\(e\,s(p(e))^{-1}\) defines \(\rho:J_1\to N/[E,N]\), written
additively in the target. It sends every \([nx]-[x]\) to the class
of \(n\), so it is independent of the representatives. Likewise
\(J_2\) is generated by \([na\mid mb]-[a\mid b]\), with \(n,m\in N\):
choose a reference lift for each nonzero base pair and a degenerate
reference lift for a pair mapping to zero. On applying \(d\), the
three differences have \(\rho\)-values equal to the class of \(m\),
minus the class of \(n(ama^{-1})\), and the class of \(n\),
respectively. Their sum is zero.

Conversely, the boundary of \([n\mid x]-[1\mid x]\) identifies
\([nx]-[x]\) with \([n]\) in \(H_1(J)\). The boundary of
\([n\mid m]-[1\mid1]\) gives \([nm]=[n]+[m]\). The boundary of
\([x\mid n]-[x\mid1]\) gives \([xnx^{-1}]=[n]\) by the first identity.
Thus \(n\mapsto[n]\) descends to \(N/[E,N]\) and is inverse to
\(\rho\). Both maps are natural by their formulas on differences.
The homology sequence of
\(0\to J_\bullet\to C_\bullet(E)\to C_\bullet(Q)\to0\) yields the
asserted sequence. The map from \(H_1(J)\) is induced by inclusion,
and the maps from \(H_i(E)\) to \(H_i(Q)\) are induced by \(p\).

To identify the connecting map with the transgression, filter the
bar complex of \(E\) by the number of entries projecting
nontrivially to \(Q\). In filtration degree \(i\), the bottom row
on the first page is \(C_i(Q)\): a tuple of length \(i+1\) with
one entry in \(N\) has filtration-preserving boundary the
difference of two lifts of the same base tuple. These differences
identify every pair of lifts. Indeed, a multiplication by \(N\)
can be moved through successive adjacent positions by inserting
the multiplier in \(N\), and then removed at the end of the
tuple. The bottom-row differential therefore projects to the bar
differential of \(Q\).

In filtration degree zero, the first-page term in total degree one
is \(H_1(N)\). The incoming first differential comes from
chains with precisely one nontrivial projected entry in degree
two and boundary in filtration zero, modulo boundaries.
The chain \([x\mid n]-[xnx^{-1}\mid x]\) has boundary
\([n]-[xnx^{-1}]\) in filtration zero, so these classes impose
conjugation. Conversely any such chain is in \(J_2\), and the
calculation with \(\rho\) shows that its boundary in filtration
zero maps to zero in \(N/[E,N]\). Thus the second-page target
is exactly \(N/[E,N]\).

For a base two-cycle \(z\), choose any lift \(\widetilde z\).
Write \(d\widetilde z\) as a sum of differences \([nx]-[x]\) and
add the chains \([n\mid x]-[1\mid x]\) with the same coefficients
to \(\widetilde z\). The new boundary is in filtration zero and
is the corresponding sum of \([n]\)'s. Its class is both the
connecting map of the short exact sequence and the second
differential \(d^2_{2,0}\) of the filtered bar complex. This fixes
the transgression sign for \(d[a\mid b]=[b]-[ab]+[a]\).
For the spectral-sequence presentation, see \cite[Section 6.8.3]{weibel}.
\end{proof}

\begin{theorem}[GR02-C04.06. Hopf's formula for the Schur multiplier]
\label{thm:gr02-hopf-formula}
For a free presentation \(1\to R\to F\to G\to 1\), the Schur
multiplier \(H_2(G;\mathbb Z)\) identifies naturally with
\[
       (R\cap[F,F])/[F,R].
\]
The natural identification compares different free presentations.
\end{theorem}

\begin{proof}
Fix a free basis \(S\) of \(F\). Its Cayley graph with edges
\(g\to gs\) for \(s\in S\) is a tree by the reduced-word normal
form \cite[K07.01.free\_group\_universal\_property]{d07}.
The augmented chain sequence
over \(\mathbb Z[F]\) has the form
\[
 0\to \bigoplus_{s\in S}\mathbb Z[F]\to\mathbb Z[F]\to\mathbb Z\to0 .
\]
Its boundary sends the edge \(g\to gs\) to \(gs-g\).
Every \(g-1\) is a path boundary, giving exactness at vertices.
A nonzero edge chain has finite support, and in the forest on
those edges a leaf has nonzero boundary coefficient. This proves
injectivity. Thus the sequence is a free resolution of the
trivial left coefficient module, and the Tor description
\cite[A10]{homalg} gives \(H_2(F;\mathbb Z)=0\).

Proposition~\ref{prop:gr02-homological-five-term} identifies
\(H_2(G;\mathbb Z)\) with the kernel of \(R/[F,R]\to F/[F,F]\),
which is the displayed quotient. For a map of free presentations,
the exact sequence is natural. A map between the quotient groups
can be lifted by choosing images of free generators. If two lifts
are chosen, their values differ elementwise by elements of the
target relation subgroup. Those differences are central modulo
commutators with the target free group, so both maps have the
same values on commutators in that quotient. Hence the induced map
on the displayed kernel is independent of the lift. Applied to
the identity of \(G\), naturality gives the canonical comparison
between presentations. See \cite[Theorem 6.8.8]{weibel}.
\end{proof}

\begin{example}[GR02-C04.07. Multipliers of cyclic groups and their products]
\label{ex:gr02-multiplier-cyclic-product}
For \(C_m\), take \(F\) free on one generator. As \([F,F]=1\),
Theorem~\ref{thm:gr02-hopf-formula} gives \(H_2(C_m;\mathbb Z)=0\).
For \(C_m\times C_n\) with \(m,n\geq 1\), present it with free
generators \(x,y\) and relations \(x^m,y^n,[x,y]\). Let \(R\)
be the normal closure of the relations. The quotient
\(D=F/[F,R]\) has central image of \(R\), which contains
\([D,D]\) because \(F/R\) is abelian. Its commutator subgroup is
generated by \(z=[x,y]\) in \(D\). In fact \(F'\) is normally
generated by \([x,y]\), since adjoining that relation makes
the generators commute, and its image in \(D\) is central.
The relations \(x^m,y^n\) have central image, so
\(z^m=z^n=1\), by the commutator rules for central commutators.
For example \([ab,y]=a[b,y]a^{-1}[a,y]=[b,y][a,y]\) in \(D\).

Put \(d=\gcd(m,n)\). If \(d>1\), map \(x,y\) to \(I+E_{12}\)
and \(I+E_{23}\) in the upper unitriangular \(3\times3\) matrices
over \(\mathbb Z/d\mathbb Z\), where \(E_{ij}\) is a matrix unit.
Their \(m\)-th and \(n\)-th powers, respectively, are identities.
Their commutator is \(I+E_{13}\), which is central in this matrix
group and has order \(d\). The image of \(R\) is central because
its normal generators have central images, so the map kills
\([F,R]\). Hence \(z\) has order at least \(d\), and its order
divides \(d\). When \(d=1\) the latter already gives \(z=1\).
The Hopf formula, with \(F'\subseteq R\), now gives
\[
          H_2(C_m\times C_n;\mathbb Z)\cong C_{\gcd(m,n)} .
\]
\end{example}

\section{The universal central extension}

A central extension of \(G\) is an extension with kernel
in the center of its middle group. It is universal if it
admits a unique homomorphism over \(G\) to every central
extension of \(G\). A group is perfect if it equals its
commutator subgroup.

\begin{theorem}[GR02-C04.08. Universal central extensions of perfect groups]
\label{thm:gr02-universal-central-extension}
A group \(G\) has a universal central extension if and only if
it is perfect. When \(G\) is perfect and \(F/R\cong G\) is any
free presentation, the map
\[
       [F,F]/[F,R]\longrightarrow G
\]
is the universal central extension. Its source is perfect,
and its kernel is \(H_2(G;\mathbb Z)\).
\end{theorem}

\begin{proof}
Suppose \(G\) is perfect, and write \(D=F/[F,R]\) and
\(\overline R\) for the image of \(R\). That image is central.
Since \([F,F]\) maps onto \(G\), the subgroup
\(U=[F,F]/[F,R]=[D,D]\) maps onto \(G\), with central kernel
\((R\cap[F,F])/[F,R]\). Theorem~\ref{thm:gr02-hopf-formula}
identifies this kernel. In addition \(D=U\overline R\), so
every commutator of \(D\) is a commutator of elements of \(U\).
It follows that \(U=[D,D]=[U,U]\).

Given a central extension \(E\to G\), lift the free generators
of \(F\) to \(E\), obtaining a homomorphism \(F\to E\) over
\(G\). Its images of the elements of \(R\) lie in the central
kernel, so the map kills \([F,R]\). Restriction yields a
homomorphism \(U\to E\) over \(G\). If two such homomorphisms
are given, their pointwise quotient lies in the central kernel
and is a homomorphism into that abelian group. It must be
trivial because \(U\) is perfect. This proves universality.

Conversely, if \(U\to G\) is universal, the split central
extension \(U/[U,U]\times G\to G\) admits two maps over \(G\)
from \(U\), by using either the zero map or the abelianization
map in the first coordinate. Uniqueness implies \(U=[U,U]\).
Surjectivity then forces \(G\) to be perfect.
See \cite[Theorem 6.9.5]{weibel}.
\end{proof}

For the stable elementary group \(E(R)\), the application in
\cite[B12]{ktheory} identifies \(K_2(R)\) with the corresponding
integral second homology.
\chapter{Finite quotients and profinite completion}
\label{ch:gr02-profinite}

A finite quotient can distinguish two elements of an infinite group.
The profinite completion puts all of these quotients into one
topological group. The compatible-tuple construction and the product
topology also occur in the Krull-topology comparison in
\emph{Advanced Field Theory} \cite[Results T04--T06]{field}. Here the
indexing objects are the finite-index normal subgroups of an arbitrary
abstract group.

\section{The inverse limit of finite quotients}

\begin{definition}[GR02-C05.01. The profinite completion]
\label{def:gr02-completion}
For an abstract group \(G\), let
\[
  \mathcal N(G)=\{N\trianglelefteq G:[G:N]<\infty\}.
\]
This set includes \(G\), which gives the trivial quotient. For
\(M,N\in\mathcal N(G)\) with \(M\subseteq N\), the transition
homomorphism \(q_{M,N}\colon G/M\to G/N\) sends \(gM\) to \(gN\).
The \emph{profinite completion} of \(G\) is
\[
  \widehat G
  =\left\{(x_N)_{N\in\mathcal N(G)}\in\prod_{N\in\mathcal N(G)}G/N
          :q_{M,N}(x_M)=x_N\text{ if }M\subseteq N\right\}.
\]
We use coordinatewise multiplication and the subspace topology of the
product of finite discrete groups. Write \(p_N\) for coordinate
projection and \(\eta_G(g)=(gN)_N\) for the canonical map.
\end{definition}

The intersection \(M\) of any finite collection of members of
\(\mathcal N(G)\) lies in \(\mathcal N(G)\), since \(G/M\) embeds into
the product of the corresponding finite quotients. In particular a
common refinement lies below each prescribed finite collection under
inclusion. The transition maps compose because each simply passes to
a larger subgroup. Thus these quotients form a cofiltered inverse
system \cite[Definition 1.2.1 and Section 1.2]{profinite-source}.

\begin{proposition}[GR02-C05.02. Topology, density, and kernel]
\label{prop:gr02-completion-topology}
The group \(\widehat G\) is compact, Hausdorff, and totally
disconnected. Each \(p_N\) is a continuous surjection, and the kernels
\(\ker p_N\) form a neighborhood basis at the identity consisting of
clopen normal subgroups. The image of \(\eta_G\) is dense and
\[
     \ker\eta_G=\bigcap_{N\in\mathcal N(G)}N .
\]
\end{proposition}

\begin{proof}
The product of the finite discrete groups is a compact Hausdorff
topological group. Each compatibility condition defines a closed
subgroup, as it is the preimage of the diagonal in
\((G/N)\times(G/N)\) under a continuous group homomorphism. Their
intersection is therefore a compact Hausdorff topological group.
Two distinct tuples have different values in some finite discrete
coordinate, so the clopen fibers of that projection separate them.
This proves total disconnectedness.

Every \(p_N\) is surjective because \(p_N\eta_G\) is the quotient
map. Basic neighborhoods of the identity constrain finitely many
coordinates to be \(1\); the kernel of \(p_M\), for the intersection
\(M\) of those subgroups, lies in the given neighborhood. Each
\(\ker p_N\) is clopen and normal. To approximate any
\(x\in\widehat G\) in a neighborhood prescribing \(x_{N_1},\ldots,x_{N_t}\),
take \(M=N_1\cap\cdots\cap N_t\) and choose \(g\in G\) with
\(gM=x_M\). Compatibility gives \(p_{N_i}\eta_G(g)=x_{N_i}\) for all
\(i\), so every such neighborhood meets \(\eta_G(G)\).
The asserted kernel follows from the formula for \(\eta_G\).
\end{proof}

For a finite \(G\), the coordinate at \(N=\{1\}\) determines all
other coordinates. Thus \(\eta_G\) is an isomorphism onto the
completion, with \(G\) given the discrete topology.

\begin{theorem}[GR02-C05.03. Universal property]
\label{thm:gr02-completion-universal}
Let \(P\) be a profinite group, regarded as a topological inverse
limit of finite discrete groups. Each abstract group homomorphism
\(f\colon G\to P\) has a unique continuous extension
\(\widehat f\colon\widehat G\to P\) satisfying
\(\widehat f\eta_G=f\).
\end{theorem}

\begin{proof}
Write \(P=\varprojlim_i P_i\) with projections \(\pi_i\). For each
\(i\), the kernel \(N_i=\ker(\pi_i f)\) has finite index in \(G\),
since \(G/N_i\) is isomorphic to a subgroup of \(P_i\).
The induced homomorphism \(G/N_i\to P_i\) composed with \(p_{N_i}\)
gives a continuous map \(h_i\colon\widehat G\to P_i\) extending
\(\pi_i f\). For each transition map in the system of the \(P_i\),
the two required coordinate maps agree on \(\eta_G(G)\), by the
compatibility of the coordinates of \(f(g)\). Their equalizer is
closed because the codomain is Hausdorff. Proposition~%
\ref{prop:gr02-completion-topology} therefore implies that they
agree everywhere. The maps \(h_i\) consequently define a
homomorphism to \(P\), continuous in the inverse-limit topology.
On \(\eta_G(G)\) its coordinates are those of \(f\).
Two continuous extensions into the Hausdorff group \(P\) agree on
the same dense subset and are equal.
\end{proof}

Taking \(P=\widehat H\) and composing a group homomorphism
\(G\to H\) with \(\eta_H\) gives its induced continuous homomorphism
on completions. Identities and composition are preserved by
the uniqueness in Theorem~\ref{thm:gr02-completion-universal}.
See \cite[Proposition 3.1.2]{profinite-source} for this comparison.

\begin{proposition}[GR02-C05.04. A cofinal family of finite quotients]
\label{prop:gr02-completion-cofinal}
If \(\mathcal D\subseteq\mathcal N(G)\) has the property that each
\(N\in\mathcal N(G)\) contains some \(D\in\mathcal D\), restriction
induces a topological group isomorphism
\[
       \widehat G \longrightarrow
       \varprojlim_{D\in\mathcal D}G/D .
\]
\end{proposition}

\begin{proof}
For \(D_1,D_2\in\mathcal D\), the condition applied to
\(D_1\cap D_2\) gives a member of \(\mathcal D\) contained in both.
Given a compatible tuple on \(\mathcal D\), define its \(N\)-coordinate
as \(q_{D,N}(x_D)\) from any \(D\subseteq N\) in \(\mathcal D\).
Two choices have a common refinement in \(\mathcal D\), which
shows that this coordinate is independent of the choice.
If \(M\subseteq N\), choosing \(D\subseteq M\) computes both
coordinates and proves their compatibility. Each coordinate of
this extension is a continuous group homomorphism since it is
given by one projection followed by a finite quotient map.
The extension and restriction recover the original coordinates
in both directions.
\end{proof}

\begin{example}[GR02-C05.05. The integers and finite fields]
\label{ex:gr02-completion-integers}
The finite-index subgroups of \(\mathbb Z\) are \(n\mathbb Z\) for
positive integers \(n\). The coordinates of the completion are
therefore the compatible residue classes in
\(\varprojlim_n\mathbb Z/n\mathbb Z\), with transition from \(m\)
to \(n\) when \(n\mid m\). One obtains the same limit using
the quotients by \(j!\mathbb Z\), since every positive integer
divides some factorial. Proposition~\ref{prop:gr02-completion-cofinal}
identifies the two topologies as well as the groups.

The absolute Galois group of \(\mathbb F_q\) is topologically
isomorphic to this completion, with \(x\mapsto x^q\) corresponding
to \(\eta_{\mathbb Z}(1)\). This is the residue comparison proved
in \emph{Advanced Field Theory} \cite[Result E06]{field}.
\end{example}

\section{Residual finiteness and three examples}

\begin{proposition}[GR02-C05.06. Residual finiteness and injectivity]
\label{prop:gr02-residual-finiteness}
An abstract group \(G\) is \emph{residually finite} if for each
\(g\ne 1\) some homomorphism to a finite group has nonidentity
value on \(g\). This is equivalent to \(\eta_G\) being injective.
In a residually finite group every finite subset maps injectively
to a single finite quotient.
\end{proposition}

\begin{proof}
A homomorphism to a finite group factors through the finite quotient
by its kernel. Conversely each quotient by \(N\in\mathcal N(G)\)
provides such a homomorphism. Thus residual finiteness is equivalent
to the intersection in Proposition~\ref{prop:gr02-completion-topology}
being trivial. Given a finite subset, choose a finite quotient
that detects \(x^{-1}y\) for each pair of distinct elements \(x,y\)
in it. The intersection of the finitely many kernels gives one
finite quotient in which all these pairs remain distinct.
For a subset of size at most one the trivial quotient suffices.
\end{proof}

\begin{proposition}[GR02-C05.07. Finitely generated abelian groups]
\label{prop:gr02-abelian-residual}
Every finitely generated abelian group is residually finite.
\end{proposition}

\begin{proof}
By the structure theorem \cite[K07.07.fg\_abelian\_extension]{d07},
such a group is isomorphic to \(\mathbb Z^r\oplus A\) with \(A\)
finite abelian. For a nonzero element with nonzero \(A\)-component,
projection to \(A\) detects it. Otherwise at least one integer
coordinate \(a\) is nonzero. Project to that coordinate and reduce
modulo \(|a|+1\), so that the residue of \(a\) is nonzero.
\end{proof}

\begin{theorem}[GR02-C05.08. Free groups]
\label{thm:gr02-free-residual}
A free group on any set of generators is residually finite.
\end{theorem}

\begin{proof}
By the reduced-word description of the free group
\cite[Result A01]{geomgroup}, a nonidentity element has a nonempty
reduced word \(w=s_{\ell}^{\varepsilon_{\ell}}\cdots
s_1^{\varepsilon_1}\), where each \(\varepsilon_i\) is \(1\) or
\(-1\) and each \(s_i\) is a generator. Use permutations acting
on the left, so the rightmost letter acts first. On
\(X=\{0,\ldots,\ell\}\), prescribe part of the permutation
\(\sigma_s\) for each generator \(s\) that occurs in \(w\):
\[
   \sigma_s(i-1)=i\quad\text{if }s_i=s,\ \varepsilon_i=1,\qquad
   \sigma_s(i)=i-1\quad\text{if }s_i=s,\ \varepsilon_i=-1.
\]
Two different steps for \(s\) could prescribe the same domain
only at a shared vertex of consecutive steps, with signs
\(-1,1\) in step order. Two different steps could prescribe the
same range only at a shared vertex with signs \(1,-1\).
Either would put adjacent inverse letters in \(w\).
Thus each prescription is a bijection between two subsets of \(X\).
Their complements have equal finite cardinalities, so choose a
bijection between the complements to extend it to a permutation.
For generators absent from the word, take the identity.

By the free-group universal property these assignments define a
homomorphism to the finite permutation group of \(X\).
For each \(i\), the letter \(s_i^{\varepsilon_i}\) moves
\(i-1\) to \(i\) by construction. The image of \(w\) consequently
moves \(0\) to \(\ell\ne 0\). This finite quotient (taking the
image of the homomorphism if necessary) detects \(w\).
The construction uses only the letters present in the word.
\end{proof}

Completing these partial permutations amounts to completing the
labelled path to a finite covering graph, as in
\cite[Theorem 3.2.1, first proof]{profinite-source}.

\begin{example}[GR02-C05.09. The additive rationals]
\label{ex:gr02-rationals-completion}
If \(f\colon(\mathbb Q,+)\to H\) has finite target of order \(n\),
then for every \(x\in\mathbb Q\)
\[
   f(x)=f(x/n)^n=1.
\]
The last equality follows from Lagrange's theorem in \(H\).
All the finite quotients are therefore trivial and so is
\(\widehat{\mathbb Q}\). Its canonical map still has dense image,
as in Proposition~\ref{prop:gr02-completion-topology}.
\end{example}

\backmatter
{\raggedright
\begin{thebibliography}{99}

\bibitem{d07}
MathGenome, \emph{Abstract Algebra} (D07), revised lecture notes,
archived in MathGenome v0.22.0. Entries have keys beginning with
\texttt{K07}.

\bibitem{homalg}
MathGenome, \emph{Advanced Homological Algebra}, Chapter~6,
Results A02, A03, A05, and A10. Companion lecture notes.

\bibitem{field}
MathGenome, \emph{Advanced Field Theory}, Chapters~2--4.
Companion lecture notes.

\bibitem{geomgroup}
MathGenome, \emph{Geometric Group Theory}, Chapters~1--3 and~6.
Companion lecture notes.

\bibitem{ktheory}
MathGenome, \emph{Algebraic K-Theory}, Chapter~2, Result B12.
Companion lecture notes.

\bibitem{chief-source}
K. Conrad, \href{https://kconrad.math.uconn.edu/blurbs/grouptheory/subgpseries2.pdf}{\emph{Subgroup Series II}},
Section~5.

\bibitem{conrad-extensions}
B. Conrad, \href{https://math.stanford.edu/~conrad/210BPage/handouts/gpext.pdf}{\emph{Group Cohomology and Group Extensions}},
Math~210B, Stanford University.

\bibitem{lassueur}
C. Lassueur, \href{https://classueur.github.io/maths/teaching/TUKL/SS21/COHOMSS21/CGSS21/COHOM_SS21.pdf}{\emph{Cohomology of Groups}},
Summer 2021, Chapters~7 and~9.

\bibitem{weibel}
C. A. Weibel, \href{https://www.cambridge.org/core/books/an-introduction-to-homological-algebra/AAA3F16482097015CD12D4376D505282}{\emph{An Introduction to Homological Algebra}},
Cambridge University Press, 1994, Sections~6.8--6.9.

\bibitem{profinite-source}
G. Wilkes, \href{https://www.dpmms.cam.ac.uk/~grw46/LectureNotes2021.pdf}{\emph{Profinite Groups and Group Cohomology}},
Cambridge Part III, Lent 2021, Chapters~1--3.

\end{thebibliography}
}
\end{document}

